% =========================================================================
\section{Literature Review}
\label{sec:literature}
% =========================================================================

Existing research is reviewed in five thematic areas that together define the problem \mahros{} addresses, followed by a comparative analysis.

\subsection{Inter-Hospital Transfer and Clinical Urgency}
Clinical studies establish that escalation transfer is frequent, structured and time-critical. Iwashyna \emph{et al.}\ mapped 47,820 Medicare critical-care transfers among 3,308 hospitals and found a stable core--periphery network in which more central hospitals had more capabilities \cite{iwashyna2009networks}. That is the hub-and-spoke topology of the Houston region studied here, where the Texas Medical Center is surrounded by community hospitals. Churpek \emph{et al.}\ associated each hour of delayed ICU transfer with about a 3\% increase in adjusted odds of mortality \cite{churpek2016delay}, which motivates time to definitive care as the primary objective. Wong and Harris showed that inter-hospital critical-care transfer delays in 49 UK hospitals were driven by organisational rather than geographical factors \cite{wong2015delays}. Mueller \emph{et al.}\ related transfer processes to patient outcomes \cite{mueller2019transfer}. Greenwood-Ericksen \emph{et al.}\ provided the most direct evidence on refusals: in 26,020 requests, acceptance fell as the receiver's census and boarding rose, and overall acceptance dropped from above 70\% before the pandemic to about 40\% afterwards \cite{greenwood2025boarding}. In the United States, the Emergency Medical Treatment and Labor Act already treats refusing a transfer despite available specialised capacity as unacceptable for emergency patients (42 CFR 489.24(f)) \cite{emtala489}. It does not cover patients who are already admitted, so it is a precedent for the accountability \mahros{} provides, not a mechanism for it.

\subsection{Centralised Coordination and Optimisation}
Operational transfer centres coordinate placements centrally. The Arizona Surge Line coordinated COVID-19 transfers across public, private and federal hospitals statewide and reported equity of access as an outcome \cite{villarroel2023surgeline}. Optimisation models go further. Dijkstra \emph{et al.}\ balanced COVID-19 patients across Dutch hospitals and regions using occupancy forecasts and a queueing network with explicit fairness \cite{dijkstra2023balancing}, and Baas \emph{et al.}\ extended that line to stochastic programming of bed capacity over demand scenarios \cite{baas2023stochastic}. Parker \emph{et al.}\ formulated linear and mixed-integer programmes, with robust variants, to redistribute demand and resources \cite{parker2020redistribution}. Mehrotra \emph{et al.}\ optimised the sharing of a ventilator stockpile among states \cite{mehrotra2020ventilator}. These approaches achieve high-quality joint allocations, including batched assignment by the Hungarian method \cite{kuhn1955hungarian}, but all of them assume that the coordinator sees each hospital's true state. \mahros{} adopts the strongest form of this approach, with optimisers that read the simulator's true private state, as its efficiency ceiling.

\subsection{Queueing and Simulation Models of Hospital Capacity}
Hospital bed pools are classically modelled as loss systems, in which a patient who arrives to a full unit is turned away or diverted. De Bruin \emph{et al.}\ validated the Erlang loss model against several years of Dutch hospital data and built a decision-support tool for ward sizing. They showed that a single target occupancy across units is unrealistic, because small units must run emptier to keep refusal rates acceptable \cite{debruin2010erlang}. Green made the same point for U.S.\ hospitals: the widely used 85\% occupancy target produces long delays in small units \cite{green2002beds}. Little's law \cite{little1961proof} links occupancy, arrival rate and length of stay in any stable system and underlies every calibration from census data. Discrete-event simulation is the standard method for evaluating hospital operations under such models \cite{gunal2010des}. \mahros{} inverts the Erlang loss relation per bed pool to recover the arrival rates that reproduce each real hospital's reported occupancy, and then validates the result from a cold start.

\subsection{Multi-Agent Negotiation and the Contract-Net Protocol}
Multi-agent systems are well established in healthcare for scheduling, decision support and monitoring \cite{isern2010agents,isern2016review}, and discrete-event simulation is the standard tool for evaluating hospital operations \cite{gunal2010des}. The Contract-Net Protocol \cite{smith1980contract}, standardised by FIPA \cite{fipa2002cnp}, allocates a task by announcement, parallel bidding and award. The manager needs no knowledge of the contractors' internal state, which makes the protocol a natural fit for transfer brokering. Automated negotiation studies the protocols and decision models of self-interested agents more broadly \cite{jennings2001negotiation}. Mechanism design seeks truthful reporting through incentives, from the sealed-bid second-price auction \cite{vickrey1961counterspeculation} to computational mechanisms \cite{nisan2001amd}. Such mechanisms need payments or priced valuations, which public-interest bed allocation does not have. \mahros{} therefore retains Contract Net's parallelism and privacy and replaces incentive compatibility with verification.

\subsection{Argumentation and Formal Dispute Resolution}
Dung's abstract argumentation frameworks \cite{dung1995acceptability} model a dispute as a graph of arguments and attacks. Of their semantics, the grounded semantics is unique, sceptical and computable in polynomial time \cite{baroni2011semantics,modgil2009algorithms}, and its labelling form and the reinstatement principle are analysed by Caminada \cite{caminada2006reinstatement}. Argumentation-based negotiation lets agents justify and criticise offers, not only make and reject them \cite{rahwan2003abn}, and Amgoud \emph{et al.}\ characterised when arguments change negotiation outcomes \cite{amgoud2007unified}. In medicine, argumentation supports clinical decision making \cite{fox2007argumentation}, and the ProCLAIM model mediated deliberation over safety-critical actions in organ transplantation \cite{tolchinsky2012deliberation}. Python solvers such as PyArg are available \cite{odekerken2022pyarg}. \mahros{} applies grounded semantics as an accountable resolution rule for evidence-backed disputes over refusals, so an outcome never depends on persuasion.

\subsection{Trust, Reputation and Enforcement}
The standard multi-agent answer to unreliable peers is trust and reputation. Agents learn whom to rely on from direct experience or third-party ratings, and surveys classify these models by information source, aggregation and robustness to manipulation \cite{ramchurn2004trust,sabater2005trust,josang2007survey}. Reputation works when interactions repeat and their outcomes are eventually observed. In transfer brokering the decisive outcome, whether the refuser actually had a bed, is never observed by anyone but the refuser, so the feedback loop reputation depends on is missing. Two other strands are more directly applicable. Economic analyses of enforcement show that random inspection combined with a credible sanction deters misconduct even at low inspection rates \cite{becker1968crime}, which motivates the random spot checks in \mahros{}. Empirical industrial organisation has detected bid rigging in procurement auctions from the joint pattern of bids \cite{porter1993bidrigging}, which motivates the co-refusal statistic used to monitor collusion.

\subsection{Privacy-Preserving Multi-Agent Coordination}
Distributed constraint optimisation (DCOP) is the formal setting closest to transfer brokering: agents with private constraints jointly choose an assignment. Léauté and Faltings showed how cryptographic techniques and distributed computation can protect four kinds of private information, namely constraint feasibility, decisions, agent identities and problem topology, with formally proven guarantees \cite{leaute2013privacy}. Greenstadt \emph{et al.}\ measured empirically how much private information standard DCOP algorithms leak \cite{greenstadt2006privacy}. These approaches assume honest inputs and carry significant computational and engineering cost. \mahros{} takes a lighter, measurable route: no occupancy value is ever transmitted, and the information each message does reveal is quantified as an entropy reduction (Section~\ref{sec:privacy}).

\subsection{Tamper-Evident Records and Digital Signatures}
Hash chains and Merkle trees make an append-only record tamper-evident \cite{merkle1988signature,haber1991timestamp}. Schneier and Kelsey extended this to audit logs on machines that may later be compromised \cite{schneier1999audit}, and Crosby and Wallach introduced history trees that give logarithmic-size membership and consistency proofs \cite{crosby2009tamper}. Blockchain designs for health records add distributed consensus over who may write \cite{azaria2016medrec}. A consortium of identifiable hospitals needs tamper \emph{evidence} rather than permissionless consensus, which \mahros{} provides with a Merkle-linked ledger and a trigger-maintained audit chain with external checkpoints. For signed refusals, Ed25519 provides fast, compact, non-repudiable signatures with 32-byte public keys and 64-byte signatures \cite{bernstein2012ed25519}.

\subsection{Comparative Analysis of Existing Approaches}
Table~\ref{tab:compare} compares \mahros{} with representative approaches along the architectural axes that matter for transfer brokering.

\begin{table*}[!t]
\centering
\caption{Comparative Analysis of Inter-Hospital Transfer Coordination Approaches}
\label{tab:compare}
\footnotesize
\setlength{\tabcolsep}{3pt}
\begin{tabularx}{\textwidth}{@{}>{\raggedright\arraybackslash}p{2.6cm}*{6}{>{\raggedright\arraybackslash}X}@{}}
\toprule
Feature & Telephone search \cite{wong2015delays} & Transfer centre / registry \cite{villarroel2023surgeline} & Central optimiser \cite{dijkstra2023balancing,parker2020redistribution} & Contract Net \cite{smith1980contract} & Argumentation-based negotiation \cite{rahwan2003abn,amgoud2007unified} & \textbf{\mahros{} (this work)} \\
\midrule
Topology & Serial, peer-to-peer & Centralised & Centralised & Decentralised & Decentralised & Decentralised multi-agent \\
Search & One hospital at a time & Parallel lookup & Joint solve & Parallel bids & Parallel, iterative & Parallel bids, 2 rounds \\
Occupancy disclosed & Verbally, per call & Every hospital, continuously & Every hospital, continuously & No & Partially & \textbf{Never} \\
Refusal verified & No & No & No (assumes truth) & No & Arguments, no evidence base & \textbf{Registry + ledger evidence} \\
Dispute resolution & None & None & None & None & Persuasion / utility & \textbf{Grounded semantics} \\
Misreporting modelled & -- & -- & No & No & Partly & \textbf{Yes, swept} \\
Collusion defence & -- & -- & -- & -- & -- & \textbf{Random spot checks} \\
Audit trail & Notes & Database & -- & -- & -- & \textbf{Hash-chained ledger + audit log} \\
Evaluation data & Clinical cohorts & Operational & Real/retrospective & Synthetic & Synthetic & \textbf{57 real hospitals} \\
\bottomrule
\end{tabularx}
\end{table*}
